\chapter{System Analysis}

\section{Existing System}
Many existing customer retention workflows in FinTech organizations are based on descriptive dashboards, periodic reports, or manually defined business rules. Common examples include flagging users with no transactions for a fixed number of days or generating monthly inactivity summaries. These approaches are useful for monitoring but do not provide predictive risk estimation.

\subsection{Disadvantages of Existing System}
\begin{itemize}
    \item No customer-level predictive scoring is available.
    \item Rule-based methods are reactive and may miss hidden nonlinear patterns.
    \item Business rules require constant manual updating.
    \item Explanations are often shallow and not tied to statistically learned patterns.
    \item It is difficult to prioritize outreach efficiently without risk stratification.
\end{itemize}

\section{Proposed System}
The proposed system is a machine learning based churn prediction framework that ingests customer data, preprocesses it, engineers features, trains predictive models, and produces a churn risk category for each customer. The system is intended to support retention teams through dashboard-based decision support.

\subsection{Advantages of Proposed System}
\begin{itemize}
    \item Predictive scoring identifies at-risk customers before churn occurs.
    \item Multiple models can be compared and tuned systematically.
    \item Risk can be categorized as High, Medium, or Low for operational use.
    \item Key features influencing churn can be highlighted to improve interpretability.
    \item The framework is extensible to future real-time or API-based deployment.
\end{itemize}

\section{Dataset Specification}
The dataset used in this report is assumed to be an anonymized FinTech customer dataset with approximately 10{,}000 records. The source is unspecified and is treated as a project assumption. The schema is inspired by public churn datasets and prior FinTech churn studies \cite{malhotra2017}.

\begin{table}[H]
\centering
\caption{Dataset Attributes}
\begin{tabular}{|p{3.2cm}|p{2.3cm}|p{6.2cm}|}
\hline
\textbf{Attribute} & \textbf{Type} & \textbf{Description} \\ \hline
CustomerID & String & Unique customer identifier removed before modeling \\ \hline
Age & Numeric & Customer age in years \\ \hline
Gender & Categorical & Gender category \\ \hline
Region & Categorical & Geographic region or service zone \\ \hline
AccountBalance & Numeric & Current account balance \\ \hline
CreditScore & Numeric & Credit score or equivalent rating \\ \hline
NumTransactions & Numeric & Number of transactions in the recent period \\ \hline
TenureMonths & Numeric & Account age in months \\ \hline
Complaints & Numeric & Number of recorded complaints \\ \hline
ActiveStatus & Categorical & Active or inactive account status \\ \hline
Subscription & Categorical & Plan type such as Basic or Gold \\ \hline
Churn & Binary & Target label where 1 indicates churn \\ \hline
\end{tabular}
\label{tab:dataset_attributes}
\end{table}

\section{Data Preprocessing}
The preprocessing pipeline includes the following steps:
\begin{itemize}
    \item removal of non-predictive identifiers such as \texttt{CustomerID};
    \item treatment of missing values using mean, median, or placeholder categories;
    \item outlier control for features such as account balance;
    \item one-hot encoding of categorical variables;
    \item scaling of numerical features when required by the model;
    \item imbalance handling using SMOTE or class weighting.
\end{itemize}

\section{Feature Engineering}
To improve predictive quality, the proposed system creates additional derived variables such as average transactions per month, complaint rate, and tenure buckets. Interaction-oriented features may also be considered when they improve validation performance. Feature importance rankings from tree-based models guide selection of the most informative predictors.

\section{System Requirements}
\subsection{Hardware Requirements}
\begin{itemize}
    \item Processor: Intel i5 or equivalent and above
    \item RAM: 8 GB minimum, 16 GB preferred
    \item Storage: 20 GB free disk space
    \item Internet connectivity for scholarly and package access
\end{itemize}

\subsection{Software Requirements}
\begin{itemize}
    \item Operating System: Windows, Linux, or macOS
    \item Programming Language: Python
    \item Libraries: pandas, NumPy, scikit-learn, XGBoost, matplotlib, seaborn
    \item Notebook or IDE: Jupyter Notebook, VS Code, or PyCharm
    \item Web Framework: Flask or Django for dashboard integration
    \item Database: PostgreSQL, MySQL, or MongoDB
\end{itemize}
